\subsection{Forms, conversion, arithmetic}\label{sub:cxforms}
\D{the two forms}{
$z=a+ib=re^{i\varphi}=r(\cos\varphi+i\sin\varphi)$; \ $\bar z=a-ib$, \ $|z|=r=\sqrt{a^2+b^2}$, \ $z\bar z=|z|^2$.\\
$\operatorname{Re}z=\frac{z+\bar z}2$, \ $\operatorname{Im}z=\frac{z-\bar z}{2i}$.
\hd{Euler:} $e^{i\varphi}=\cos\varphi+i\sin\varphi$, \ $e^{i\pi}=-1$, \ $e^{2\pi i}=1$.
}
\R{convert both ways --- the quadrant table settles it}{
\hd{Polar \too\ Cartesian} (never ambiguous): $a=r\cos\varphi$, \ $b=r\sin\varphi$.\\
\hd{Cartesian \too\ polar:} $r=\sqrt{a^2+b^2}$, and read $\varphi$ off the sign of $a,b$ --- \textbf{do not think, look it up}:
\begin{tsTabular}{@{}l@{\quad}l@{\quad}l@{}}
\hd{$a>0$} (right half) & $\varphi=\arctan\frac ba$ & quadrant I / IV\\
\hd{$a<0$, $b\ge0$} & $\varphi=\arctan\frac ba+\pi$ & quadrant II\\
\hd{$a<0$, $b<0$} & $\varphi=\arctan\frac ba-\pi$ & quadrant III\\
\hd{$a=0$} & $\varphi=+\frac\pi2$ if $b>0$, $-\frac\pi2$ if $b<0$ & $\varphi$ undefined at $z=0$
\end{tsTabular}
\hd{Why:} $\arctan$ only ever returns a value in $(-\frac\pi2,\frac\pi2)$, i.e.\ it always answers in the \emph{right} half-plane; for $a<0$ you must rotate by $\pi$ to reach the correct half. The \textbf{sign} of that $\pi$ is chosen so the answer stays in the principal range $(-\pi,\pi]$ --- $+\pi$ would overshoot past $\pi$ when $b<0$. \warn A $\varphi$ outside $(-\pi,\pi]$ is not wrong, just not principal: add or subtract $2\pi$ to bring it back.\\
\hd{Ex.} $z=-1-i$: $r=\sqrt2$, $\arctan\frac{-1}{-1}=\frac\pi4$, third quadrant \ra\ $\varphi=\frac\pi4-\pi=-\frac{3\pi}4$. \ ($+\pi$ would give $\frac{5\pi}4>\pi$.)
}
\F{arithmetic rules}{
\hd{Multiply/divide in polar:} $r_1e^{i\varphi_1}\cdot r_2e^{i\varphi_2}=r_1r_2e^{i(\varphi_1+\varphi_2)}$, \ $\frac{z}{w}=\frac{r_z}{r_w}e^{i(\varphi_z-\varphi_w)}$.\\
\hd{Divide in Cartesian:} multiply by $\frac{\bar w}{\bar w}$: $\frac zw=\frac{z\bar w}{|w|^2}$.\\
\hd{Conjugate:} $\overline{z+w}=\bar z+\bar w$, $\overline{zw}=\bar z\bar w$, $\overline{z/w}=\bar z/\bar w$, $\bar{\bar z}=z$, $z\in\Rl\iff z=\bar z$, $|\bar z|=|z|$.\\
\hd{de Moivre:} $(\cos\varphi+i\sin\varphi)^n=\cos(n\varphi)+i\sin(n\varphi)$.
}
\R{use Euler to derive trig identities on the spot}{
$\cos x=\frac{e^{ix}+e^{-ix}}2$, $\sin x=\frac{e^{ix}-e^{-ix}}{2i}$. Multiply the exponentials, collect. \\
\hd{Ex.} $\cos^2x=\left(\frac{e^{ix}+e^{-ix}}2\right)^2=\frac{e^{2ix}+2+e^{-2ix}}4=\frac{1+\cos2x}2$. Same trick gives every double/half-angle formula, so you never have to memorise them.
}

\subsection{Roots and equations}\label{sub:cxroots}
\R{solve $z^n=w$ ($n$-th roots)}{
\begin{rcp}
\item Write $w=\rho e^{i\psi}$.
\item $z_k=\rho^{1/n}e^{i(\psi+2\pi k)/n}$ for $k=0,1,\dots,n-1$.
\item Exactly $n$ roots, all on the circle of radius $\rho^{1/n}$, spaced by $\frac{2\pi}n$.
\end{rcp}
\hd{Roots of unity:} $z^n=1$ \ra\ $z_k=e^{2\pi ik/n}$; their sum is $0$ for $n\ge2$ and their product is $(-1)^{n+1}$.
}
\R{quadratic with complex data}{
$az^2+bz+c=0\Rightarrow z=\frac{-b\pm\sqrt{D}}{2a}$, $D=b^2-4ac$.
\begin{bul}
\item Real coefficients, $D<0$: $z=-\frac b{2a}\pm i\frac{\sqrt{|D|}}{2a}$ (a conjugate pair).
\item Complex $D$: get $\sqrt D$ with the $n$-th-root recipe at $n=2$ (two roots, take either).
\end{bul}
}
\F{complex log and general powers $+$ DEF: what a \emph{branch} is}{
\hd{The problem:} $e^{i\varphi}=e^{i(\varphi+2\pi k)}$, so $\arg z$ is only fixed \textbf{mod $2\pi$} and ``$\log$'' has \emph{infinitely many} values $\log z=\ln|z|+i(\arg z+2\pi k)$, $k\in\Z$. That is not a function.\\
\hd{A branch} $=$ one consistent choice of $k$. The \hd{principal branch} is the choice $\operatorname{Arg}z\in(-\pi,\pi]$ (the same range as the quadrant table above), giving the genuine function
\[\operatorname{Log}z=\ln|z|+i\operatorname{Arg}z,\qquad z^w:=e^{w\operatorname{Log}z}.\]
\hd{Ex.} $\operatorname{Log}(-1)=\ln1+i\pi=i\pi$. \ $i^i=e^{i\operatorname{Log}i}=e^{i\cdot i\pi/2}=e^{-\pi/2}$ --- \emph{real}.\\
\warn Only $\operatorname{Log}$ is a function; $\sqrt{zw}=\sqrt z\sqrt w$ and $\log(zw)=\log z+\log w$ hold for the multivalued $\log$ but can be off by $2\pi i$ once you pin down the principal branch. Counterexample $z=w=-1$: $\sqrt{zw}=1$ but $\sqrt z\sqrt w=i\cdot i=-1$.
}

\subsection{Polynomials: what is guaranteed}\label{sub:poly}
\F{Fundamental theorem of algebra}{
$p$ of degree $n\ge1$ with \emph{any} complex coefficients has \textbf{exactly $n$ roots in $\Cx$ counted with multiplicity}, and factors as $p(z)=a_n\prod_{j}(z-z_j)^{m_j}$ with $\sum m_j=n$.
}
\X{the three classic multiple-choice errors}{
\begin{bul}
\item \claimF{``$n$ roots, all \emph{distinct}''} --- only \emph{with multiplicity} ($z^5$ has one root of multiplicity 5).
\item \claimF{``complex roots come in conjugate pairs'' \emph{in general}} --- \yes\ \textbf{only if all coefficients are real}. Degree 5 with complex coefficients need not have any real root.
\item \claimT{``odd degree \ra\ a real root''} --- \textbf{only for real coefficients} (IVT on $p$, which changes sign at $\pm\infty$).
\end{bul}
}
\R{factor a polynomial over $\Rl$ / over $\Cx$}{
\begin{rcp}
\item Guess a root among the divisors of the constant term (over $\Z$), or spot $z=\pm1$.
\item Polynomial division by $(z-z_0)$, repeat on the quotient.
\item Over $\Cx$: keep going until all factors are linear.
\item Over $\Rl$: pair each non-real root with its conjugate to get an irreducible quadratic $z^2-2\operatorname{Re}(z_0)z+|z_0|^2$. Real factorisation = linear factors $\times$ irreducible quadratics only.
\end{rcp}
}
